% GENERATED by docs/manual/build.py - do not edit; change the source it names instead.
% source: scripts/watchover.py (TOOLS, TOOL_HOW, TOOL_EXITS) and each command's argparse parser
\chapter{Kit commands}\label{ch:kit}
Run in the governance repository (or \texttt{template/} of an installed skill): \texttt{scripts/watchover <command> [options]}; \texttt{<command> --help} or \texttt{help <command>} prints the entry in the terminal. \emph{Live} commands call Akeyless; the others are offline. Each entry ends with a real execution (Section~\ref{sec:examples}).

\section{\texttt{test}}\label{kit:test}
\begin{woentry}
\refrow{Synopsis}{\texttt{scripts/\allowbreak{}watchover test [options]}\hfill\emph{offline}}
\refrow{Does}{Run every offline check (what CI runs).}
\refrow{How}{\begin{itemize}[nosep,leftmargin=1.2em]
\item Runs, stopping at the first failure: fmt -check, init -backend=false, validate and \texttt{test} (mocked provider) in terraform/ and bootstrap/; guard\_messages.sh (each negative fixture fails with the intended guard); export\_roundtrip.sh; \texttt{simulate}; pytest.
\end{itemize}
}
\refrow{Exit codes}{0 all passed; otherwise the exit code of the first failing step}
\end{woentry}
{\small
\begin{longtable}{@{}>{\raggedright\arraybackslash}p{0.27\linewidth}>{\raggedright\arraybackslash}p{0.17\linewidth}>{\raggedright\arraybackslash}p{0.13\linewidth}>{\raggedright\arraybackslash}p{0.35\linewidth}@{}}
\toprule
\textbf{Option} & \textbf{Value} & \textbf{Required / default} & \textbf{Description} \\ \midrule\endhead
\texttt{\mbox{-\allowbreak{}}\mbox{-\allowbreak{}}tool} & \texttt{terraform\textbar{}\allowbreak{}tofu} & default terraform & IaC binary to use \\
\texttt{\mbox{-\allowbreak{}}\mbox{-\allowbreak{}}skip-\allowbreak{}iac} & -- & optional & only run the Python checks \\
\bottomrule
\end{longtable}}

\paragraph{Examples}\leavevmode
\begin{Verbatim}[fontsize=\footnotesize,breaklines,breakanywhere]
scripts/watchover test
scripts/watchover test --tool tofu
\end{Verbatim}
\paragraph{Sample run}\leavevmode
% GENERATED by docs/manual/build.py - do not edit; change the source it names instead.
% source: docs/manual/build.py EXAMPLES (executed in a sandbox copy of the repository)
Every offline check, as CI runs them (formatting, validation, Terraform tests with a mocked provider, guard messages, export round trip, bootstrap tests, Python tests, access simulation):
\begin{terminal}
\begin{Verbatim}[fontsize=\scriptsize,breaklines,breakanywhere,commandchars=\\\{\}]
\textcolor{wo-prompt}{\$} \textbf{scripts/watchover test --tool tofu}

$ tofu fmt -check -recursive   (in terraform)

$ tofu init -backend=false -input=false   (in terraform)

Initializing provider plugins...
- terraform.io/builtin/terraform is built in to OpenTofu
- Reusing previous version of akeyless-community/akeyless from the dependency lock file
- Installing akeyless-community/akeyless v2.0.5...
- Installed akeyless-community/akeyless v2.0.5 (unauthenticated)

OpenTofu has been successfully initialized!

$ tofu validate   (in terraform)
Success! The configuration is valid.

$ tofu test   (in terraform)
tests/model.tftest.hcl... pass
  run "reference_model_plans"... pass
  run "gateway_objects_need_gateway_api"... pass
  run "ara_rule_allowed_when_enabled"... pass
  run "root_scope_allowed_with_flag"... pass
  run "dup"... pass
  run "deny_mixed"... pass
  run "bad_capability"... pass
  run "bad_rule_type"... pass
  run "sra_conflict"... pass
  run "root_scope"... pass
  run "unknown_am"... pass
  run "missing_subclaims"... pass
  run "bad_prefix"... pass
  run "bad_mgmt"... pass
  run "bad_issuer"... pass
  run "bad_gateway"... pass
  run "bad_event"... pass
  run "scaffolded_forwarders_plan"... pass
  run "bad_periodic"... pass
  run "slack_no_url"... pass
  run "empty_gateway_api_is_unset"... pass
  run "naming_check_disabled_with_empty_string"... pass
[... 37 lines omitted ...]
Success! The configuration is valid.

$ tofu test   (in bootstrap)
tests/bootstrap.tftest.hcl... pass
  run "github"... pass
  run "gitlab"... pass
  run "bad_platform"... pass

Success! 3 passed, 0 failed.

$ /usr/local/bin/python3 ./scripts/akl_simulate.py --model terraform/model --expectations terraform/model/access_expectations.yaml   (in .)
PASS  payments-admin           item        /payments/checkout/nonprod/db/app                  expect=['create', 'delete', 'list', 'read', 'update'] got=['create', 'delete', 'list', 'read', 'update']
PASS  payments-admin           item        /payments/checkout/prod/db/app                     expect=['list', 'read'] got=['list', 'read']
PASS  payments-admin           item        /payments/checkout/prod/break-glass/root           expect=[] got=[]
PASS  payments-admin           item        /hr/payroll/prod/db                                expect=[] got=[]
PASS  auditor                  item        /payments/checkout/prod/db/app                     expect=['list'] got=['list']
PASS  auditor                  role        /governance/roles/payments/team-admin              expect=['list', 'read'] got=['list', 'read']
PASS  random-employee          item        /payments/checkout/prod/db/app                     expect=[] got=[]
PASS  checkout-prod-pod        item        /payments/checkout/prod/db/app                     expect=['list', 'read'] got=['list', 'read']
PASS  checkout-prod-pod        item        /payments/checkout/prod/break-glass/root           expect=[] got=[]
PASS  checkout-prod-pod        item        /payments/checkout/nonprod/db/app                  expect=[] got=[]
PASS  checkout-prod-pod        item        /pki/issuers/payments-server                       expect=['list', 'read'] got=['list', 'read']

0 mismatch(es)

$ /usr/local/bin/python3 -m pytest -q scripts/tests   (in .)
.......................................................................                      [100%]
71 passed in <t>s

all checks passed
\end{Verbatim}
\exitcode{0}
\end{terminal}

Only the Python checks:
\begin{terminal}
\begin{Verbatim}[fontsize=\scriptsize,breaklines,breakanywhere,commandchars=\\\{\}]
\textcolor{wo-prompt}{\$} \textbf{scripts/watchover test --skip-iac}

$ /usr/local/bin/python3 ./scripts/akl_simulate.py --model terraform/model --expectations terraform/model/access_expectations.yaml   (in .)
PASS  payments-admin           item        /payments/checkout/nonprod/db/app                  expect=['create', 'delete', 'list', 'read', 'update'] got=['create', 'delete', 'list', 'read', 'update']
PASS  payments-admin           item        /payments/checkout/prod/db/app                     expect=['list', 'read'] got=['list', 'read']
PASS  payments-admin           item        /payments/checkout/prod/break-glass/root           expect=[] got=[]
PASS  payments-admin           item        /hr/payroll/prod/db                                expect=[] got=[]
PASS  auditor                  item        /payments/checkout/prod/db/app                     expect=['list'] got=['list']
PASS  auditor                  role        /governance/roles/payments/team-admin              expect=['list', 'read'] got=['list', 'read']
PASS  random-employee          item        /payments/checkout/prod/db/app                     expect=[] got=[]
PASS  checkout-prod-pod        item        /payments/checkout/prod/db/app                     expect=['list', 'read'] got=['list', 'read']
PASS  checkout-prod-pod        item        /payments/checkout/prod/break-glass/root           expect=[] got=[]
PASS  checkout-prod-pod        item        /payments/checkout/nonprod/db/app                  expect=[] got=[]
PASS  checkout-prod-pod        item        /pki/issuers/payments-server                       expect=['list', 'read'] got=['list', 'read']

0 mismatch(es)

$ /usr/local/bin/python3 -m pytest -q scripts/tests   (in .)
.......................................................................                      [100%]
71 passed in <t>s

all checks passed
\end{Verbatim}
\exitcode{0}
\end{terminal}



\section{\texttt{simulate}}\label{kit:simulate}
\begin{woentry}
\refrow{Synopsis}{\texttt{scripts/\allowbreak{}watchover simulate [options]}\hfill\emph{offline}}
\refrow{Does}{Offline: evaluate access\_expectations.yaml against the model.}
\refrow{How}{\begin{itemize}[nosep,leftmargin=1.2em]
\item Loads model/roles/*.yaml and auth\_methods.yaml.
\item For each check: selects the roles whose associations match the persona's auth method and claims; collects the rules of the right type whose path pattern matches (claim templates expanded); a matching deny removes everything.
\item Compares the resulting capability set with \texttt{expect} (exact match).
\end{itemize}
}
\refrow{Exit codes}{0 all met; 1 mismatch; 2 model error}
\end{woentry}
{\small
\begin{longtable}{@{}>{\raggedright\arraybackslash}p{0.27\linewidth}>{\raggedright\arraybackslash}p{0.17\linewidth}>{\raggedright\arraybackslash}p{0.13\linewidth}>{\raggedright\arraybackslash}p{0.35\linewidth}@{}}
\toprule
\textbf{Option} & \textbf{Value} & \textbf{Required / default} & \textbf{Description} \\ \midrule\endhead
\texttt{\mbox{-\allowbreak{}}\mbox{-\allowbreak{}}model} & \texttt{MODEL} & \textbf{required} & model directory (terraform/model) \\
\texttt{\mbox{-\allowbreak{}}\mbox{-\allowbreak{}}expectations} & \texttt{EXPECTATIONS} & \textbf{required} & access\_expectations.yaml \\
\texttt{\mbox{-\allowbreak{}}\mbox{-\allowbreak{}}explain} & -- & optional & print the matching rules for every check, not only failures \\
\bottomrule
\end{longtable}}

\paragraph{Examples}\leavevmode
\begin{Verbatim}[fontsize=\footnotesize,breaklines,breakanywhere]
scripts/watchover simulate --model terraform/model --expectations terraform/model/access_expectations.yaml --explain
\end{Verbatim}
\paragraph{Sample run}\leavevmode
% GENERATED by docs/manual/build.py - do not edit; change the source it names instead.
% source: docs/manual/build.py EXAMPLES (executed in a sandbox copy of the repository)
Offline evaluation of the access expectations against the model:
\begin{terminal}
\begin{Verbatim}[fontsize=\scriptsize,breaklines,breakanywhere,commandchars=\\\{\}]
\textcolor{wo-prompt}{\$} \textbf{scripts/watchover simulate --model terraform/model --expectations terraform/model/access_expectations.yaml}
PASS  payments-admin           item        /payments/checkout/nonprod/db/app                  expect=['create', 'delete', 'list', 'read', 'update'] got=['create', 'delete', 'list', 'read', 'update']
PASS  payments-admin           item        /payments/checkout/prod/db/app                     expect=['list', 'read'] got=['list', 'read']
PASS  payments-admin           item        /payments/checkout/prod/break-glass/root           expect=[] got=[]
PASS  payments-admin           item        /hr/payroll/prod/db                                expect=[] got=[]
PASS  auditor                  item        /payments/checkout/prod/db/app                     expect=['list'] got=['list']
PASS  auditor                  role        /governance/roles/payments/team-admin              expect=['list', 'read'] got=['list', 'read']
PASS  random-employee          item        /payments/checkout/prod/db/app                     expect=[] got=[]
PASS  checkout-prod-pod        item        /payments/checkout/prod/db/app                     expect=['list', 'read'] got=['list', 'read']
PASS  checkout-prod-pod        item        /payments/checkout/prod/break-glass/root           expect=[] got=[]
PASS  checkout-prod-pod        item        /payments/checkout/nonprod/db/app                  expect=[] got=[]
PASS  checkout-prod-pod        item        /pki/issuers/payments-server                       expect=['list', 'read'] got=['list', 'read']

0 mismatch(es)
\end{Verbatim}
\exitcode{0}
\end{terminal}

With \texttt{-\allowbreak{}{}-\allowbreak{}explain}, every check shows the rules that produced the result:
\begin{terminal}
\begin{Verbatim}[fontsize=\scriptsize,breaklines,breakanywhere,commandchars=\\\{\}]
\textcolor{wo-prompt}{\$} \textbf{scripts/watchover simulate --model terraform/model --expectations terraform/model/access_expectations.yaml --explain | head -n 12}
PASS  payments-admin           item        /payments/checkout/nonprod/db/app                  expect=['create', 'delete', 'list', 'read', 'update'] got=['create', 'delete', 'list', 'read', 'update']
        /governance/roles/payments/team-admin: /payments/+/nonprod/* -> ['create', 'delete', 'list', 'read', 'update']
PASS  payments-admin           item        /payments/checkout/prod/db/app                     expect=['list', 'read'] got=['list', 'read']
        /governance/roles/payments/team-admin: /payments/+/prod/* -> ['list', 'read']
PASS  payments-admin           item        /payments/checkout/prod/break-glass/root           expect=[] got=[]
        /governance/roles/payments/team-admin: /payments/+/prod/* -> ['list', 'read']
        /governance/roles/payments/team-admin: /payments/+/prod/break-glass/* -> ['deny']
        deny overrides all other permissions
PASS  payments-admin           item        /hr/payroll/prod/db                                expect=[] got=[]
        no matching rule
PASS  auditor                  item        /payments/checkout/prod/db/app                     expect=['list'] got=['list']
        /governance/roles/platform/auditors: /* -> ['list']
\end{Verbatim}
\exitcode{0}
\end{terminal}



\section{\texttt{verify}}\label{kit:verify}
\begin{woentry}
\refrow{Synopsis}{\texttt{scripts/\allowbreak{}watchover verify [options]}\hfill\emph{live}}
\refrow{Does}{Live: compare expectations with describe-permissions run as each persona's CLI profile.}
\refrow{How}{\begin{itemize}[nosep,leftmargin=1.2em]
\item For each check whose persona has a \texttt{profile}: runs \texttt{akeyless describe-\allowbreak{}permissions \mbox{-\allowbreak{}}\mbox{-\allowbreak{}}path \mbox{-\allowbreak{}}\mbox{-\allowbreak{}}type \mbox{-\allowbreak{}}\mbox{-\allowbreak{}}json \mbox{-\allowbreak{}}\mbox{-\allowbreak{}}profile \textless{}persona\textgreater{}} and compares \texttt{client\_\allowbreak{}permissions} with \texttt{expect}.
\item Personas without a profile are skipped and counted.
\end{itemize}
}
\refrow{Exit codes}{0 all met; 1 mismatch; 2 CLI call failed}
\end{woentry}
{\small
\begin{longtable}{@{}>{\raggedright\arraybackslash}p{0.27\linewidth}>{\raggedright\arraybackslash}p{0.17\linewidth}>{\raggedright\arraybackslash}p{0.13\linewidth}>{\raggedright\arraybackslash}p{0.35\linewidth}@{}}
\toprule
\textbf{Option} & \textbf{Value} & \textbf{Required / default} & \textbf{Description} \\ \midrule\endhead
\texttt{\mbox{-\allowbreak{}}\mbox{-\allowbreak{}}expectations} & \texttt{EXPECTATIONS} & \textbf{required} & access\_expectations.yaml (personas need a 'profile') \\
\texttt{\mbox{-\allowbreak{}}\mbox{-\allowbreak{}}akeyless} & \texttt{AKEYLESS} & default akeyless & path to the akeyless CLI binary \\
\texttt{\mbox{-\allowbreak{}}\mbox{-\allowbreak{}}error-\allowbreak{}means-\allowbreak{}no-\allowbreak{}access} & -- & optional & treat a failing describe-permissions call as 'no permissions' (check the CLI behaviour first) \\
\bottomrule
\end{longtable}}

\paragraph{Examples}\leavevmode
\begin{Verbatim}[fontsize=\footnotesize,breaklines,breakanywhere]
scripts/watchover verify --expectations terraform/model/access_expectations.yaml
\end{Verbatim}
\paragraph{Sample run}\leavevmode
% GENERATED by docs/manual/build.py - do not edit; change the source it names instead.
% source: docs/manual/build.py EXAMPLES (executed in a sandbox copy of the repository)
The stand-in account has one console change that is not in Git:
\begin{terminal}
\begin{Verbatim}[fontsize=\scriptsize,breaklines,breakanywhere,commandchars=\\\{\}]
\textcolor{wo-prompt}{\$} \textbf{scripts/watchover verify --expectations terraform/model/access_expectations.yaml}
PASS  payments-admin           /payments/checkout/nonprod/db/app                  expect=['create', 'delete', 'list', 'read', 'update'] got=['create', 'delete', 'list', 'read', 'update']
PASS  payments-admin           /payments/checkout/prod/db/app                     expect=['list', 'read'] got=['list', 'read']
PASS  payments-admin           /payments/checkout/prod/break-glass/root           expect=[] got=[]
PASS  payments-admin           /hr/payroll/prod/db                                expect=[] got=[]
PASS  auditor                  /payments/checkout/prod/db/app                     expect=['list'] got=['list']
[... 1 lines omitted ...]
FAIL  random-employee          /payments/checkout/prod/db/app                     expect=[] got=['list', 'read']
PASS  checkout-prod-pod        /payments/checkout/prod/db/app                     expect=['list', 'read'] got=['list', 'read']
PASS  checkout-prod-pod        /payments/checkout/prod/break-glass/root           expect=[] got=[]
PASS  checkout-prod-pod        /payments/checkout/nonprod/db/app                  expect=[] got=[]
PASS  checkout-prod-pod        /pki/issuers/payments-server                       expect=['list', 'read'] got=['list', 'read']

1 mismatch(es), 0 check(s) skipped (persona without profile)
\end{Verbatim}
\exitcode{1}
\end{terminal}



\section{\texttt{export}}\label{kit:export}
\begin{woentry}
\refrow{Synopsis}{\texttt{scripts/\allowbreak{}watchover export [options]}\hfill\emph{live}}
\refrow{Does}{Export existing config and translate it: YAML model + import blocks, or plain Terraform (\mbox{-}\mbox{-}hcl).}
\refrow{How}{\begin{itemize}[nosep,leftmargin=1.2em]
\item Reads the selected object families through the CLI (\texttt{\mbox{-\allowbreak{}}\mbox{-\allowbreak{}}json}) or the REST API (token), following pagination; filters by name, root folder, depth or persona.
\item Writes raw/ (JSON, secrets redacted) and REPORT.md (items needing a decision).
\item Model mode: model/*.imported.yaml + imports.tf for the kit. \mbox{-}\mbox{-}hcl: hcl/imports.tf with one import block per object; \mbox{-}\mbox{-}generate runs \texttt{plan -\allowbreak{}generate-\allowbreak{}config-\allowbreak{}out} so Terraform writes the resources.
\end{itemize}
}
\refrow{Exit codes}{0 written; 2 read/HTTP error or invalid options}
\end{woentry}
{\small
\begin{longtable}{@{}>{\raggedright\arraybackslash}p{0.27\linewidth}>{\raggedright\arraybackslash}p{0.17\linewidth}>{\raggedright\arraybackslash}p{0.13\linewidth}>{\raggedright\arraybackslash}p{0.35\linewidth}@{}}
\toprule
\textbf{Option} & \textbf{Value} & \textbf{Required / default} & \textbf{Description} \\ \midrule\endhead
\texttt{\mbox{-\allowbreak{}}\mbox{-\allowbreak{}}out} & \texttt{OUT} & \textbf{required} & output directory (model/, imports.tf or hcl/, raw/, REPORT.md) \\
\texttt{\mbox{-\allowbreak{}}\mbox{-\allowbreak{}}transport} & \texttt{cli\textbar{}\allowbreak{}api} & default cli & how to read Akeyless: the akeyless CLI (profiles, SSO logins) or the REST API directly (token) \\
\texttt{\mbox{-\allowbreak{}}\mbox{-\allowbreak{}}profile} & \texttt{PROFILE} & optional & cli transport: akeyless CLI profile with read access to the objects to export \\
\texttt{\mbox{-\allowbreak{}}\mbox{-\allowbreak{}}akeyless} & \texttt{AKEYLESS} & default akeyless & cli transport: path to the akeyless CLI binary \\
\texttt{\mbox{-\allowbreak{}}\mbox{-\allowbreak{}}gateway-\allowbreak{}url} & \texttt{GATEWAY\_\allowbreak{}URL} & optional & cli transport, \mbox{-}\mbox{-}types gateway: Gateway configuration URL (https://gw:8000) \\
\texttt{\mbox{-\allowbreak{}}\mbox{-\allowbreak{}}api} & \texttt{API} & default https://api.akeyless.io & api transport: API base URL (default \$AKEYLESS\_API or https://api.akeyless.io) \\
\texttt{\mbox{-\allowbreak{}}\mbox{-\allowbreak{}}token} & \texttt{TOKEN} & optional & api transport: Akeyless token (default \$AKEYLESS\_TOKEN; e.g. from 'watchover token') \\
\texttt{\mbox{-\allowbreak{}}\mbox{-\allowbreak{}}gateway-\allowbreak{}api} & \texttt{GATEWAY\_\allowbreak{}API} & optional & api transport, \mbox{-}\mbox{-}types gateway: Gateway V2 API (https://gw:8000/api/v2) \\
\texttt{\mbox{-\allowbreak{}}\mbox{-\allowbreak{}}from-\allowbreak{}dir} & \texttt{FROM\_\allowbreak{}DIR} & optional & replay a raw/ directory of a previous export instead of calling the CLI \\
\texttt{\mbox{-\allowbreak{}}\mbox{-\allowbreak{}}types} & \texttt{TYPES} & default roles,auth-methods,folders,custom-fields,key-policies,pki-issuers,account-settings,targets,groups & comma-separated object families (default: all but gateway): roles, auth-methods, folders, custom-fields, key-policies, pki-issuers, account-settings, targets, groups, gateway \\
\texttt{\mbox{-\allowbreak{}}\mbox{-\allowbreak{}}only} & \texttt{ONLY} & optional & keep only objects whose name matches this regex \\
\texttt{\mbox{-\allowbreak{}}\mbox{-\allowbreak{}}exclude} & \texttt{EXCLUDE} & \emph{repeatable} & regex of names to skip (repeatable) \\
\texttt{\mbox{-\allowbreak{}}\mbox{-\allowbreak{}}root} & \texttt{ROOT} & default / & folder to start the item/folder walk from \\
\texttt{\mbox{-\allowbreak{}}\mbox{-\allowbreak{}}max-\allowbreak{}depth} & \texttt{MAX\_\allowbreak{}DEPTH} & default 20 & maximum folder depth of the walk \\
\texttt{\mbox{-\allowbreak{}}\mbox{-\allowbreak{}}persona-\allowbreak{}auth-\allowbreak{}method} & \texttt{PERSONA\_\allowbreak{}AUTH\_\allowbreak{}METHOD} & optional & keep one identity's view: its auth method ... \\
\texttt{\mbox{-\allowbreak{}}\mbox{-\allowbreak{}}persona-\allowbreak{}claim} & \texttt{PERSONA\_\allowbreak{}CLAIM} & \emph{repeatable} & ... and its claims, key=value (repeatable); keeps roles whose association matches \\
\texttt{\mbox{-\allowbreak{}}\mbox{-\allowbreak{}}hcl} & -- & optional & write plain Terraform import blocks for every object (hcl/) \\
\texttt{\mbox{-\allowbreak{}}\mbox{-\allowbreak{}}generate} & -- & optional & with \mbox{-}\mbox{-}hcl: run 'terraform plan -generate-config-out' (needs terraform and AKEYLESS\_ACCESS\_ID/AKEYLESS\_ACCESS\_KEY in the environment) \\
\texttt{\mbox{-\allowbreak{}}\mbox{-\allowbreak{}}terraform} & \texttt{TERRAFORM} & default terraform & terraform or tofu binary used by \mbox{-}\mbox{-}generate \\
\bottomrule
\end{longtable}}

\paragraph{Examples}\leavevmode
\begin{Verbatim}[fontsize=\footnotesize,breaklines,breakanywhere]
scripts/watchover export --out export/ --profile admin --exclude '^/governance/auth/ci-terraform$'
scripts/watchover export --out export/ --profile admin --types auth-methods --only sso
scripts/watchover export --out export/ --profile admin --persona-auth-method /governance/auth/sso-entra-oidc --persona-claim email=alice@example.com --persona-claim groups=grp-akl-payments-admins
AKEYLESS_ACCESS_ID=p-... AKEYLESS_ACCESS_KEY=... scripts/watchover export --out export/ --profile admin --hcl --generate
AKEYLESS_TOKEN=$(scripts/watchover token) scripts/watchover export --out export/ --transport api --hcl --types targets,groups,gateway --gateway-api https://gw.example.com:8000/api/v2
scripts/watchover export --out export2/ --from-dir export/raw     # offline replay
\end{Verbatim}
\paragraph{Sample run}\leavevmode
% GENERATED by docs/manual/build.py - do not edit; change the source it names instead.
% source: docs/manual/build.py EXAMPLES (executed in a sandbox copy of the repository)
\begin{terminal}
\begin{Verbatim}[fontsize=\scriptsize,breaklines,breakanywhere,commandchars=\\\{\}]
\textcolor{wo-prompt}{\$} \textbf{scripts/watchover export --out export/ --profile admin --types roles,auth-methods && head -n 8 export/imports.tf}
exported (model): 3 roles, 7 auth-methods; 7 item(s) in REPORT.md
# Generated by watchover export (model mode) - addresses of terraform/*.tf.
# Remove after the first successful apply.

import \{
  to = akeyless_auth_method_aws_iam.this["/governance/auth/aws-payments-prod"]
  id = "/governance/auth/aws-payments-prod"
\}
\end{Verbatim}
\exitcode{0}
\end{terminal}



\section{\texttt{reconcile}}\label{kit:reconcile}
\begin{woentry}
\refrow{Synopsis}{\texttt{scripts/\allowbreak{}watchover reconcile [options]}\hfill\emph{live}}
\refrow{Does}{Compare the governance repository (e.g. GitHub Enterprise) with live Akeyless; write a report and a patch to review. Read-only: never commits or pushes.}
\refrow{How}{\begin{itemize}[nosep,leftmargin=1.2em]
\item Gets the repository read-only: shallow clone of \mbox{-}\mbox{-}ref with the push URL disabled, or reads \mbox{-}\mbox{-}repo-dir as it is. Never writes to it, never commits or pushes.
\item Exports roles, auth methods, folders, custom fields and key policies from Akeyless (as \texttt{export}; or uses \mbox{-}\mbox{-}export-dir / \mbox{-}\mbox{-}from-dir).
\item Compares object by object after normalisation (defaults, ordering, *\_var secret keys and guard switches ignored): changed, only in Akeyless, only in the repository.
\item Writes REPORT.md, changes.patch (\texttt{git apply}), proposed/ (complete changed files at their repository paths, edited in place so comments survive), imports.tf for objects added from Akeyless, summary.json and the export used.
\end{itemize}
}
\refrow{Exit codes}{0 no differences; 1 differences found (files written); 2 git, export or model error}
\refrow{Environment}{\texttt{AKEYLESS\_\allowbreak{}TOKEN}: api transport: token when \mbox{-}\mbox{-}token is not given\newline \texttt{AKEYLESS\_\allowbreak{}API}: api transport: API URL}
\end{woentry}
{\small
\begin{longtable}{@{}>{\raggedright\arraybackslash}p{0.27\linewidth}>{\raggedright\arraybackslash}p{0.17\linewidth}>{\raggedright\arraybackslash}p{0.13\linewidth}>{\raggedright\arraybackslash}p{0.35\linewidth}@{}}
\toprule
\textbf{Option} & \textbf{Value} & \textbf{Required / default} & \textbf{Description} \\ \midrule\endhead
\texttt{\mbox{-\allowbreak{}}\mbox{-\allowbreak{}}repo-\allowbreak{}url} & \texttt{REPO\_\allowbreak{}URL} & optional & governance repository to clone read-only (GitHub Enterprise HTTPS or SSH URL; uses your git credentials, e.g. \texttt{gh auth login \mbox{-\allowbreak{}}\mbox{-\allowbreak{}}hostname \textless{}ghe-\allowbreak{}host\textgreater{}}) \\
\texttt{\mbox{-\allowbreak{}}\mbox{-\allowbreak{}}repo-\allowbreak{}dir} & \texttt{REPO\_\allowbreak{}DIR} & optional & existing local checkout to compare (read only, never modified) \\
\texttt{\mbox{-\allowbreak{}}\mbox{-\allowbreak{}}ref} & \texttt{REF} & default main & branch to compare with \mbox{-}\mbox{-}repo-url \\
\texttt{\mbox{-\allowbreak{}}\mbox{-\allowbreak{}}model-\allowbreak{}path} & \texttt{MODEL\_\allowbreak{}PATH} & default terraform/model & model directory inside the repository \\
\texttt{\mbox{-\allowbreak{}}\mbox{-\allowbreak{}}out} & \texttt{OUT} & default reconcile & output directory (replaced; keep it out of Git) \\
\texttt{\mbox{-\allowbreak{}}\mbox{-\allowbreak{}}export-\allowbreak{}dir} & \texttt{EXPORT\_\allowbreak{}DIR} & optional & use an earlier \texttt{watchover export} output instead of exporting now \\
\texttt{\mbox{-\allowbreak{}}\mbox{-\allowbreak{}}from-\allowbreak{}dir} & \texttt{FROM\_\allowbreak{}DIR} & optional & export offline from a raw/ directory (as \texttt{export \mbox{-\allowbreak{}}\mbox{-\allowbreak{}}from-\allowbreak{}dir}) \\
\texttt{\mbox{-\allowbreak{}}\mbox{-\allowbreak{}}transport} & \texttt{cli\textbar{}\allowbreak{}api} & default cli & how to read Akeyless (as \texttt{export}) \\
\texttt{\mbox{-\allowbreak{}}\mbox{-\allowbreak{}}profile} & \texttt{PROFILE} & optional & cli transport: akeyless CLI profile (a read-only identity is enough) \\
\texttt{\mbox{-\allowbreak{}}\mbox{-\allowbreak{}}akeyless} & \texttt{AKEYLESS} & default akeyless & cli transport: path to the akeyless CLI binary \\
\texttt{\mbox{-\allowbreak{}}\mbox{-\allowbreak{}}api} & \texttt{API} & optional & api transport: API URL (default \$AKEYLESS\_API or https://api.akeyless.io) \\
\texttt{\mbox{-\allowbreak{}}\mbox{-\allowbreak{}}token} & \texttt{TOKEN} & optional & api transport: Akeyless token (default \$AKEYLESS\_TOKEN) \\
\texttt{\mbox{-\allowbreak{}}\mbox{-\allowbreak{}}only} & \texttt{ONLY} & optional & compare only objects whose name matches this regex \\
\texttt{\mbox{-\allowbreak{}}\mbox{-\allowbreak{}}exclude} & \texttt{EXCLUDE} & \emph{repeatable} & regex of names to skip (repeatable) \\
\texttt{\mbox{-\allowbreak{}}\mbox{-\allowbreak{}}include-\allowbreak{}unmodelled} & -- & optional & also propose fields set in Akeyless but absent from the model \\
\texttt{\mbox{-\allowbreak{}}\mbox{-\allowbreak{}}propose-\allowbreak{}removals} & -- & optional & propose removing objects that are in the repository but not in Akeyless \\
\bottomrule
\end{longtable}}

\paragraph{Examples}\leavevmode
\begin{Verbatim}[fontsize=\footnotesize,breaklines,breakanywhere]
scripts/watchover reconcile --repo-url https://ghe.example.com/platform/akeyless-governance.git --profile akl-readonly
scripts/watchover reconcile --repo-dir ~/src/akeyless-governance --ref main --transport api
scripts/watchover reconcile --repo-dir . --export-dir export/ --out /tmp/reconcile
scripts/watchover reconcile --repo-url git@ghe.example.com:platform/akeyless-governance.git --profile ro --only '^/governance/roles/payments/' --propose-removals
\end{Verbatim}
\paragraph{Sample run}\leavevmode
% GENERATED by docs/manual/build.py - do not edit; change the source it names instead.
% source: docs/manual/build.py EXAMPLES (executed in a sandbox copy of the repository)
The stand-in GitHub Enterprise repository has one role changed in the console, one role created there and one role only in Git: (Shown as you would run it; executed against the local stand-in.)
\begin{terminal}
\begin{Verbatim}[fontsize=\scriptsize,breaklines,breakanywhere,commandchars=\\\{\}]
\textcolor{wo-prompt}{\$} \textbf{scripts/watchover reconcile --repo-url https://ghe.example.com/platform/akeyless-governance.git --profile akl-readonly --out reconcile/; sed -n '20,31p' reconcile/REPORT.md}
exported (model): 3 roles, 7 auth-methods, 1 folders, 1 custom-fields, 1 key-policies; 8 item(s) in REPORT.md
reconcile: 1 changed, 1 only in Akeyless, 1 only in the repository; 2 file(s) in the proposal
review reconcile/REPORT.md and reconcile/changes.patch - nothing was committed or pushed
## Changed outside Git (1)

### roles `/governance/roles/payments/team-admin`  - `terraform/model/roles/payments.yaml`
- **rules** changed:
  - only in the repository: `\{"capabilities": ["list", "read"], "path": "/payments/+/nonprod/*"\}`
  - only in Akeyless: `\{"capabilities": ["create", "list", "read"], "path": "/payments/+/nonprod/*"\}`

## Only in Akeyless (1)

### roles `/governance/roles/platform/god`  - `terraform/model/roles/imported.yaml`
- created outside Git: proposed as an addition + import block in `imports.tf`
\end{Verbatim}
\exitcode{1}
\end{terminal}



\section{\texttt{forwarder}}\label{kit:forwarder}
\begin{woentry}
\refrow{Synopsis}{\texttt{scripts/\allowbreak{}watchover forwarder [options]}\hfill\emph{offline}}
\refrow{Does}{Scaffold a new event forwarder (model entry or HCL); \mbox{-}\mbox{-}import-existing adopts one by name.}
\refrow{How}{\begin{itemize}[nosep,leftmargin=1.2em]
\item Validates name, type, event types, sources and runner with the same rules as the Terraform guards.
\item Builds the model entry; secrets become \texttt{*\_\allowbreak{}var} keys of TF\_VAR\_event\_forwarder\_secrets.
\item Prints YAML (default) or HCL; \mbox{-}\mbox{-}write inserts the entry into governance.yaml without touching the rest of the file and refuses duplicates; \mbox{-}\mbox{-}import-existing adds an import block (ID = name).
\item Presets: certificates = certificate-pending-expiration, certificate-expired, certificate-provisioning-success, certificate-provisioning-failure; auth-methods = auth-method-pending-expiration, auth-method-expired, multi-auth-failure, uid-rotation-failure; rotation = next-automatic-rotation, rotated-secret-success, rotated-secret-failure, rotation-usage; secrets = dynamic-secret-failure, static-secret-updated, secret-sync; access-requests = request-access, apply-justification, email-auth-method-approved; operations = gateway-inactive, rate-limiting, usage, usage-report; critical = certificate-expired, certificate-provisioning-failure, auth-method-expired, rotated-secret-failure, dynamic-secret-failure, multi-auth-failure, uid-rotation-failure, gateway-inactive.
\end{itemize}
}
\refrow{Exit codes}{0 written; 2 invalid input or duplicate (nothing written)}
\end{woentry}
{\small
\begin{longtable}{@{}>{\raggedright\arraybackslash}p{0.27\linewidth}>{\raggedright\arraybackslash}p{0.17\linewidth}>{\raggedright\arraybackslash}p{0.13\linewidth}>{\raggedright\arraybackslash}p{0.35\linewidth}@{}}
\toprule
\textbf{Option} & \textbf{Value} & \textbf{Required / default} & \textbf{Description} \\ \midrule\endhead
\texttt{\mbox{-\allowbreak{}}\mbox{-\allowbreak{}}name} & \texttt{NAME} & \textbf{required} & forwarder name, e.g. /governance/events/secops-mail \\
\texttt{\mbox{-\allowbreak{}}\mbox{-\allowbreak{}}type} & \texttt{email\textbar{}\allowbreak{}webhook\textbar{}\allowbreak{}slack\textbar{}\allowbreak{}teams\textbar{}\allowbreak{}service\_\allowbreak{}now} & \textbf{required} & forwarder type \\
\texttt{\mbox{-\allowbreak{}}\mbox{-\allowbreak{}}description} & \texttt{DESCRIPTION} & optional & description \\
\texttt{\mbox{-\allowbreak{}}\mbox{-\allowbreak{}}preset} & \texttt{access-\allowbreak{}requests\textbar{}\allowbreak{}auth-\allowbreak{}methods\textbar{}\allowbreak{}certificates\textbar{}\allowbreak{}critical\textbar{}\allowbreak{}operations\textbar{}\allowbreak{}rotation\textbar{}\allowbreak{}secrets} & \emph{repeatable} & add a group of event types (repeatable); the events of each preset are listed under 'How it works' \\
\texttt{\mbox{-\allowbreak{}}\mbox{-\allowbreak{}}event} & \texttt{EVENT} & \emph{repeatable} & add one event type (repeatable) \\
\texttt{\mbox{-\allowbreak{}}\mbox{-\allowbreak{}}items-\allowbreak{}source} & \texttt{ITEMS\_\allowbreak{}SOURCE} & \emph{repeatable} & item path pattern to watch, e.g. '/pki/*' (repeatable) \\
\texttt{\mbox{-\allowbreak{}}\mbox{-\allowbreak{}}targets-\allowbreak{}source} & \texttt{TARGETS\_\allowbreak{}SOURCE} & \emph{repeatable} & target path pattern to watch (repeatable) \\
\texttt{\mbox{-\allowbreak{}}\mbox{-\allowbreak{}}auth-\allowbreak{}method-\allowbreak{}source} & \texttt{AUTH\_\allowbreak{}METHOD\_\allowbreak{}SOURCE} & \emph{repeatable} & auth method path pattern to watch (repeatable) \\
\texttt{\mbox{-\allowbreak{}}\mbox{-\allowbreak{}}gateway-\allowbreak{}source} & \texttt{GATEWAY\_\allowbreak{}SOURCE} & \emph{repeatable} & Gateway cluster URL to watch (repeatable; required for teams) \\
\texttt{\mbox{-\allowbreak{}}\mbox{-\allowbreak{}}runner} & \texttt{immediate\textbar{}\allowbreak{}periodic} & default immediate & send at once or as a periodic digest \\
\texttt{\mbox{-\allowbreak{}}\mbox{-\allowbreak{}}every} & \texttt{EVERY} & optional & hours between digests (required with \mbox{-}\mbox{-}runner periodic) \\
\texttt{\mbox{-\allowbreak{}}\mbox{-\allowbreak{}}disabled} & -- & optional & create it disabled \\
\texttt{\mbox{-\allowbreak{}}\mbox{-\allowbreak{}}key} & \texttt{KEY} & optional & protection key for the forwarder's secret (default: account key) \\
\texttt{\mbox{-\allowbreak{}}\mbox{-\allowbreak{}}email-\allowbreak{}to} & \texttt{EMAIL\_\allowbreak{}TO} & \emph{repeatable} & email: recipient (repeatable) \\
\texttt{\mbox{-\allowbreak{}}\mbox{-\allowbreak{}}include-\allowbreak{}error} & -- & optional & email: include error details \\
\texttt{\mbox{-\allowbreak{}}\mbox{-\allowbreak{}}override-\allowbreak{}url} & \texttt{OVERRIDE\_\allowbreak{}URL} & optional & email: Gateway URL used in links (port 18888) \\
\texttt{\mbox{-\allowbreak{}}\mbox{-\allowbreak{}}url} & \texttt{URL} & optional & webhook: URL (slack/teams URLs are secrets and go in the secrets map) \\
\texttt{\mbox{-\allowbreak{}}\mbox{-\allowbreak{}}auth-\allowbreak{}type} & \texttt{AUTH\_\allowbreak{}TYPE} & optional & webhook: user-pass/bearer-token/certificate; service\_now: user-pass/jwt \\
\texttt{\mbox{-\allowbreak{}}\mbox{-\allowbreak{}}username} & \texttt{USERNAME} & optional & webhook: user for user-pass \\
\texttt{\mbox{-\allowbreak{}}\mbox{-\allowbreak{}}host} & \texttt{HOST} & optional & service\_now: instance host \\
\texttt{\mbox{-\allowbreak{}}\mbox{-\allowbreak{}}admin-\allowbreak{}name} & \texttt{ADMIN\_\allowbreak{}NAME} & optional & service\_now: admin name (user-pass) \\
\texttt{\mbox{-\allowbreak{}}\mbox{-\allowbreak{}}client-\allowbreak{}id} & \texttt{CLIENT\_\allowbreak{}ID} & optional & service\_now: client ID (jwt) \\
\texttt{\mbox{-\allowbreak{}}\mbox{-\allowbreak{}}user-\allowbreak{}email} & \texttt{USER\_\allowbreak{}EMAIL} & optional & service\_now: user email (jwt) \\
\texttt{\mbox{-\allowbreak{}}\mbox{-\allowbreak{}}format} & \texttt{yaml\textbar{}\allowbreak{}hcl} & default yaml & model entry (yaml) or plain resource (hcl) \\
\texttt{\mbox{-\allowbreak{}}\mbox{-\allowbreak{}}write} & \texttt{GOVERNANCE\_\allowbreak{}YAML} & optional & yaml: insert the entry into this governance.yaml \\
\texttt{\mbox{-\allowbreak{}}\mbox{-\allowbreak{}}import-\allowbreak{}existing} & -- & optional & the forwarder already exists in Akeyless: also print the import block (import ID = name) \\
\bottomrule
\end{longtable}}

\paragraph{Examples}\leavevmode
\begin{Verbatim}[fontsize=\footnotesize,breaklines,breakanywhere]
scripts/watchover forwarder --name /governance/events/secops-mail --type email --email-to secops@example.com --preset critical --items-source '/*' --auth-method-source '/*'
scripts/watchover forwarder --name /governance/events/pki-slack --type slack --preset certificates --items-source '/pki/*' --write terraform/model/governance.yaml
scripts/watchover forwarder --name /governance/events/siem --type webhook --url https://siem.example.com/akeyless --auth-type bearer-token --preset critical --auth-method-source '/*' --format hcl
scripts/watchover forwarder --name /governance/events/legacy-mail --type email --email-to ops@example.com --event certificate-expired --items-source '/*' --import-existing
\end{Verbatim}
\paragraph{Sample run}\leavevmode
% GENERATED by docs/manual/build.py - do not edit; change the source it names instead.
% source: docs/manual/build.py EXAMPLES (executed in a sandbox copy of the repository)
A model entry (default output):
\begin{terminal}
\begin{Verbatim}[fontsize=\scriptsize,breaklines,breakanywhere,commandchars=\\\{\}]
\textcolor{wo-prompt}{\$} \textbf{scripts/watchover forwarder --name /governance/events/secops-mail --type email --email-to secops@example.com --preset critical --items-source '/*' --auth-method-source '/*'}
# add under event_forwarders: in terraform/model/governance.yaml
  /governance/events/secops-mail:
    type: email
    event_types:
    - certificate-expired
    - certificate-provisioning-failure
    - auth-method-expired
    - rotated-secret-failure
    - dynamic-secret-failure
    - multi-auth-failure
    - uid-rotation-failure
    - gateway-inactive
    runner_type: immediate
    items_event_source_locations:
    - /*
    auth_methods_event_source_locations:
    - /*
    email_to:
    - secops@example.com
\end{Verbatim}
\exitcode{0}
\end{terminal}

Insert it into the model (comments and the rest of the file are kept; the secret key is printed): (The sandbox copy of governance.yaml is restored afterwards, so later examples see the original.)
\begin{terminal}
\begin{Verbatim}[fontsize=\scriptsize,breaklines,breakanywhere,commandchars=\\\{\}]
\textcolor{wo-prompt}{\$} \textbf{scripts/watchover forwarder --name /governance/events/pki-slack --type slack --preset certificates --items-source '/pki/*' --runner periodic --every 24 --write terraform/model/governance.yaml && grep -n -A3 'pki-slack' terraform/model/governance.yaml}
added /governance/events/pki-slack to terraform/model/governance.yaml
59:  /governance/events/pki-slack:
60-    type: slack
61-    event_types:
62-    - certificate-pending-expiration

# provide these keys in TF_VAR_event_forwarder_secrets (CI secret, JSON map):
#   governance_events_pki_slack_url
\end{Verbatim}
\exitcode{0}
\end{terminal}

Plain HCL plus an import block for a forwarder that already exists:
\begin{terminal}
\begin{Verbatim}[fontsize=\scriptsize,breaklines,breakanywhere,commandchars=\\\{\}]
\textcolor{wo-prompt}{\$} \textbf{scripts/watchover forwarder --name /governance/events/siem --type webhook --url https://siem.example.com/akeyless --auth-type bearer-token --preset critical --auth-method-source '/*' --format hcl --import-existing}
resource "akeyless_event_forwarder_webhook" "governance_events_siem" \{
  name                                = "/governance/events/siem"
  event_types                         = ["certificate-expired", "certificate-provisioning-failure", "auth-method-expired", "rotated-secret-failure", "dynamic-secret-failure", "multi-auth-failure", "uid-rotation-failure", "gateway-inactive"]
  runner_type                         = "immediate"
  auth_methods_event_source_locations = ["/*"]
  url                                 = "https://siem.example.com/akeyless"
  auth_type                           = "bearer-token"
  auth_token                          = var.event_forwarder_secrets["governance_events_siem_token"]
\}

# imports.tf - the forwarder exists already; Terraform adopts it on the next apply
import \{
  to = akeyless_event_forwarder_webhook.governance_events_siem
  id = "/governance/events/siem"
\}

# needs: variable "event_forwarder_secrets" \{ type = map(string), sensitive = true \} (already declared in terraform/main.tf)

# provide these keys in TF_VAR_event_forwarder_secrets (CI secret, JSON map):
#   governance_events_siem_token
\end{Verbatim}
\exitcode{0}
\end{terminal}

ServiceNow with JWT authentication, created disabled:
\begin{terminal}
\begin{Verbatim}[fontsize=\scriptsize,breaklines,breakanywhere,commandchars=\\\{\}]
\textcolor{wo-prompt}{\$} \textbf{scripts/watchover forwarder --name /governance/events/snow --type service_now --host example.service-now.com --auth-type jwt --client-id snow-client --user-email akeyless@example.com --preset rotation --items-source '/*' --disabled}
# add under event_forwarders: in terraform/model/governance.yaml
  /governance/events/snow:
    type: service_now
    event_types:
    - next-automatic-rotation
    - rotated-secret-success
    - rotated-secret-failure
    - rotation-usage
    runner_type: immediate
    enable: false
    items_event_source_locations:
    - /*
    host: example.service-now.com
    auth_type: jwt
    client_id: snow-client
    user_email: akeyless@example.com
    client_secret_var: governance_events_snow_client_secret
    app_private_key_var: governance_events_snow_app_private_key

# provide these keys in TF_VAR_event_forwarder_secrets (CI secret, JSON map):
#   governance_events_snow_client_secret
#   governance_events_snow_app_private_key
\end{Verbatim}
\exitcode{0}
\end{terminal}

Invalid input is rejected before Terraform runs (exit code 2):
\begin{terminal}
\begin{Verbatim}[fontsize=\scriptsize,breaklines,breakanywhere,commandchars=\\\{\}]
\textcolor{wo-prompt}{\$} \textbf{scripts/watchover forwarder --name /governance/events/ops --type teams --preset operations --items-source '/*'}
error: teams needs --gateway-source (provider 2.0.5 requires gateway_event_source_locations)
\end{Verbatim}
\exitcode{2}
\end{terminal}



\section{\texttt{conjur}}\label{kit:conjur}
\begin{woentry}
\refrow{Synopsis}{\texttt{scripts/\allowbreak{}watchover conjur [options]}\hfill\emph{offline}}
\refrow{Does}{Translate Conjur policy into model roles + identity bindings.}
\refrow{How}{\begin{itemize}[nosep,leftmargin=1.2em]
\item Parses the policy files (!policy, !variable, !host, !layer, !group, !user, !permit, !grant), resolving ids against their policy namespace.
\item Maps privileges to capabilities (execute -\textgreater{} read, read -\textgreater{} list, update -\textgreater{} update) on the migrated item paths; one role per layer or group; nested grants are flattened.
\item Writes roles/conjur.yaml, bindings.yaml (host/user -\textgreater{} roles, auth method TODO) and REPORT.md.
\end{itemize}
}
\refrow{Exit codes}{0 written; 2 unreadable or invalid policy}
\end{woentry}
{\small
\begin{longtable}{@{}>{\raggedright\arraybackslash}p{0.27\linewidth}>{\raggedright\arraybackslash}p{0.17\linewidth}>{\raggedright\arraybackslash}p{0.13\linewidth}>{\raggedright\arraybackslash}p{0.35\linewidth}@{}}
\toprule
\textbf{Option} & \textbf{Value} & \textbf{Required / default} & \textbf{Description} \\ \midrule\endhead
\texttt{policy} & \texttt{POLICY...} & \textbf{required} & Conjur policy YAML files (loaded at the root namespace in the given order) \\
\texttt{\mbox{-\allowbreak{}}\mbox{-\allowbreak{}}item-\allowbreak{}prefix} & \texttt{ITEM\_\allowbreak{}PREFIX} & default /conjur & Akeyless folder the Gateway migration imports into \\
\texttt{\mbox{-\allowbreak{}}\mbox{-\allowbreak{}}role-\allowbreak{}prefix} & \texttt{ROLE\_\allowbreak{}PREFIX} & default /governance/roles/ & prefix of the generated Akeyless role names \\
\texttt{\mbox{-\allowbreak{}}\mbox{-\allowbreak{}}out} & \texttt{OUT} & \textbf{required} & output directory (roles/conjur.yaml, bindings.yaml, REPORT.md) \\
\bottomrule
\end{longtable}}

\paragraph{Examples}\leavevmode
\begin{Verbatim}[fontsize=\footnotesize,breaklines,breakanywhere]
scripts/watchover conjur root.yml prod.yml --item-prefix /conjur --out export/conjur
\end{Verbatim}
\paragraph{Sample run}\leavevmode
% GENERATED by docs/manual/build.py - do not edit; change the source it names instead.
% source: docs/manual/build.py EXAMPLES (executed in a sandbox copy of the repository)
\begin{terminal}
\begin{Verbatim}[fontsize=\scriptsize,breaklines,breakanywhere,commandchars=\\\{\}]
\textcolor{wo-prompt}{\$} \textbf{scripts/watchover conjur scripts/tests/fixtures/conjur.yml --out export/conjur && sed -n '1,9p' export/conjur/roles/conjur.yaml}
2 roles, 2 identities, 1 note(s) -> export/conjur
roles:
  /governance/roles/conjur/group-payments-admins:
    allow_external_auth_methods: true
    description: Translated from Conjur group 'payments-admins'
    rules:
    - capabilities:
      - list
      - update
      path: /conjur/prod/payments/db/password
\end{Verbatim}
\exitcode{0}
\end{terminal}



\section{\texttt{account-settings}}\label{kit:account-settings}
\begin{woentry}
\refrow{Synopsis}{\texttt{scripts/\allowbreak{}watchover account-\allowbreak{}settings [options]}\hfill\emph{live}}
\refrow{Does}{Apply account\_settings\_cli (settings the provider cannot manage) with the CLI.}
\refrow{How}{\begin{itemize}[nosep,leftmargin=1.2em]
\item Reads account\_settings\_cli from governance.yaml; rejects unknown flags and keys Terraform manages.
\item Builds one \texttt{akeyless update-\allowbreak{}account-\allowbreak{}settings} call (lists repeat the flag, booleans as =true/false); \mbox{-}\mbox{-}dry-run prints it, otherwise runs it with \mbox{-}\mbox{-}profile or \mbox{-}\mbox{-}token.
\end{itemize}
}
\refrow{Exit codes}{0 applied/printed; 2 invalid settings; else the CLI's exit code}
\refrow{Environment}{\texttt{AKEYLESS\_\allowbreak{}TOKEN}: token used when \mbox{-}\mbox{-}token is not given}
\end{woentry}
{\small
\begin{longtable}{@{}>{\raggedright\arraybackslash}p{0.27\linewidth}>{\raggedright\arraybackslash}p{0.17\linewidth}>{\raggedright\arraybackslash}p{0.13\linewidth}>{\raggedright\arraybackslash}p{0.35\linewidth}@{}}
\toprule
\textbf{Option} & \textbf{Value} & \textbf{Required / default} & \textbf{Description} \\ \midrule\endhead
\texttt{\mbox{-\allowbreak{}}\mbox{-\allowbreak{}}governance} & \texttt{GOVERNANCE} & \textbf{required} & path to model/governance.yaml (reads account\_settings\_cli) \\
\texttt{\mbox{-\allowbreak{}}\mbox{-\allowbreak{}}profile} & \texttt{PROFILE} & optional & akeyless CLI profile of an account admin \\
\texttt{\mbox{-\allowbreak{}}\mbox{-\allowbreak{}}token} & \texttt{TOKEN} & optional & Akeyless token (default: \$AKEYLESS\_TOKEN, e.g. from 'watchover token') \\
\texttt{\mbox{-\allowbreak{}}\mbox{-\allowbreak{}}akeyless} & \texttt{AKEYLESS} & default akeyless & path to the akeyless CLI binary \\
\texttt{\mbox{-\allowbreak{}}\mbox{-\allowbreak{}}dry-\allowbreak{}run} & -- & optional & print the update-account-settings command, do not run it \\
\bottomrule
\end{longtable}}

\paragraph{Examples}\leavevmode
\begin{Verbatim}[fontsize=\footnotesize,breaklines,breakanywhere]
scripts/watchover account-settings --governance terraform/model/governance.yaml --dry-run
AKEYLESS_TOKEN=$(scripts/watchover token) scripts/watchover account-settings --governance terraform/model/governance.yaml
\end{Verbatim}
\paragraph{Sample run}\leavevmode
% GENERATED by docs/manual/build.py - do not edit; change the source it names instead.
% source: docs/manual/build.py EXAMPLES (executed in a sandbox copy of the repository)
\begin{terminal}
\begin{Verbatim}[fontsize=\scriptsize,breaklines,breakanywhere,commandchars=\\\{\}]
\textcolor{wo-prompt}{\$} \textbf{scripts/watchover account-settings --governance terraform/model/governance.yaml --dry-run}
+ akeyless update-account-settings --allowed-client-type=cli --allowed-client-type=ui --allowed-client-type=sdk --bound-ips=10.0.0.0/8 --lock-allowed-client-type=true --lock-bound-ips=true --max-rotation-interval=90 --max-rotation-interval-enable=true --usage-event-enable=true --usage-event-interval=90 --usage-event-object-type=item
\end{Verbatim}
\exitcode{0}
\end{terminal}



\section{\texttt{token}}\label{kit:token}
\begin{woentry}
\refrow{Synopsis}{\texttt{scripts/\allowbreak{}watchover token [options]}\hfill\emph{live}}
\refrow{Does}{CI: exchange the pipeline JWT for an Akeyless token.}
\refrow{How}{\begin{itemize}[nosep,leftmargin=1.2em]
\item POSTs \{access-id, access-type: jwt, jwt\} to \textless{}api\textgreater{}/auth and prints the token. No credential is stored.
\end{itemize}
}
\refrow{Exit codes}{0 token printed; 1 missing settings or authentication failed}
\refrow{Environment}{\texttt{AKEYLESS\_\allowbreak{}API}: API base URL\newline \texttt{AKEYLESS\_\allowbreak{}ACCESS\_\allowbreak{}ID}: CI auth method access ID\newline \texttt{CI\_\allowbreak{}JWT}: CI-issued JWT}
\end{woentry}
{\small
\begin{longtable}{@{}>{\raggedright\arraybackslash}p{0.27\linewidth}>{\raggedright\arraybackslash}p{0.17\linewidth}>{\raggedright\arraybackslash}p{0.13\linewidth}>{\raggedright\arraybackslash}p{0.35\linewidth}@{}}
\toprule
\textbf{Option} & \textbf{Value} & \textbf{Required / default} & \textbf{Description} \\ \midrule\endhead
\texttt{\mbox{-\allowbreak{}}\mbox{-\allowbreak{}}api} & \texttt{API} & optional & Akeyless API (default: \$AKEYLESS\_API or https://api.akeyless.io) \\
\texttt{\mbox{-\allowbreak{}}\mbox{-\allowbreak{}}access-\allowbreak{}id} & \texttt{ACCESS\_\allowbreak{}ID} & optional & access ID of the CI auth method (default: \$AKEYLESS\_ACCESS\_ID) \\
\texttt{\mbox{-\allowbreak{}}\mbox{-\allowbreak{}}jwt} & \texttt{JWT} & optional & the CI-issued JWT (default: \$CI\_JWT; prefer the env var - arguments are visible in the process list) \\
\bottomrule
\end{longtable}}

\paragraph{Examples}\leavevmode
\begin{Verbatim}[fontsize=\footnotesize,breaklines,breakanywhere]
CI_JWT=$JWT AKEYLESS_ACCESS_ID=p-xxxx scripts/watchover token
\end{Verbatim}
\paragraph{Sample run}\leavevmode
% GENERATED by docs/manual/build.py - do not edit; change the source it names instead.
% source: docs/manual/build.py EXAMPLES (executed in a sandbox copy of the repository)
(Shown as you would run it; executed against the local stand-in.)
\begin{terminal}
\begin{Verbatim}[fontsize=\scriptsize,breaklines,breakanywhere,commandchars=\\\{\}]
\textcolor{wo-prompt}{\$} \textbf{CI_JWT=$CI_JWT AKEYLESS_ACCESS_ID=p-ci0demo scripts/watchover token}
t-demo-0b3f9c
\end{Verbatim}
\exitcode{0}
\end{terminal}



\section{\texttt{build-skill}}\label{kit:build-skill}
\begin{woentry}
\refrow{Synopsis}{\texttt{scripts/\allowbreak{}watchover build-\allowbreak{}skill [options]}\hfill\emph{offline}}
\refrow{Does}{Build the agent-neutral skill package (Agent Skills ZIP) and the prompt pack.}
\refrow{How}{\begin{itemize}[nosep,leftmargin=1.2em]
\item Copies SKILL.md (+ agents/ sidecar), docs/*.md as references/ and the kit as template/; validates the result; zips it.
\item Writes the prompt pack: full and short system prompt + knowledge files.
\end{itemize}
}
\refrow{Exit codes}{0 written; 2 built skill invalid}
\end{woentry}
{\small
\begin{longtable}{@{}>{\raggedright\arraybackslash}p{0.27\linewidth}>{\raggedright\arraybackslash}p{0.17\linewidth}>{\raggedright\arraybackslash}p{0.13\linewidth}>{\raggedright\arraybackslash}p{0.35\linewidth}@{}}
\toprule
\textbf{Option} & \textbf{Value} & \textbf{Required / default} & \textbf{Description} \\ \midrule\endhead
\texttt{\mbox{-\allowbreak{}}\mbox{-\allowbreak{}}out} & \texttt{OUT} & optional & output directory, replaced (default: dist/ in the kit root) \\
\bottomrule
\end{longtable}}

\paragraph{Examples}\leavevmode
\begin{Verbatim}[fontsize=\footnotesize,breaklines,breakanywhere]
scripts/watchover build-skill
scripts/watchover build-skill --out /tmp/wo-dist
\end{Verbatim}
\paragraph{Sample run}\leavevmode
% GENERATED by docs/manual/build.py - do not edit; change the source it names instead.
% source: docs/manual/build.py EXAMPLES (executed in a sandbox copy of the repository)
Package the skill for upload:
\begin{terminal}
\begin{Verbatim}[fontsize=\scriptsize,breaklines,breakanywhere,commandchars=\\\{\}]
\textcolor{wo-prompt}{\$} \textbf{scripts/watchover build-skill && unzip -l dist/akeyless-watchover.zip | tail -n 1}
built dist/akeyless-watchover.zip (156K)
   337065                     167 files
\end{Verbatim}
\exitcode{0}
\end{terminal}



\section{\texttt{install-skill}}\label{kit:install-skill}
\begin{woentry}
\refrow{Synopsis}{\texttt{scripts/\allowbreak{}watchover install-\allowbreak{}skill [options]}\hfill\emph{offline}}
\refrow{Does}{Install the skill where an AI agent looks for skills (per-agent folders or any folder).}
\refrow{How}{\begin{itemize}[nosep,leftmargin=1.2em]
\item Builds the package in a temporary folder and copies it to the agent's skills folder (project or user scope, or \mbox{-}\mbox{-}dest). Never overwrites without \mbox{-}\mbox{-}force.
\end{itemize}
}
\refrow{Exit codes}{0 installed; 2 target exists (no \mbox{-}\mbox{-}force) or custom without \mbox{-}\mbox{-}dest}
\end{woentry}
{\small
\begin{longtable}{@{}>{\raggedright\arraybackslash}p{0.27\linewidth}>{\raggedright\arraybackslash}p{0.17\linewidth}>{\raggedright\arraybackslash}p{0.13\linewidth}>{\raggedright\arraybackslash}p{0.35\linewidth}@{}}
\toprule
\textbf{Option} & \textbf{Value} & \textbf{Required / default} & \textbf{Description} \\ \midrule\endhead
\texttt{\mbox{-\allowbreak{}}\mbox{-\allowbreak{}}agent} & \texttt{agents\textbar{}\allowbreak{}claude\textbar{}\allowbreak{}codex\textbar{}\allowbreak{}copilot\textbar{}\allowbreak{}cursor\textbar{}\allowbreak{}gemini\textbar{}\allowbreak{}custom} & \textbf{required} & target agent; 'agents' = the shared .agents/skills folder; 'custom' needs \mbox{-}\mbox{-}dest \\
\texttt{\mbox{-\allowbreak{}}\mbox{-\allowbreak{}}scope} & \texttt{project\textbar{}\allowbreak{}user} & default project & project = inside \mbox{-}\mbox{-}project; user = your home directory \\
\texttt{\mbox{-\allowbreak{}}\mbox{-\allowbreak{}}project} & \texttt{PROJECT} & default . & repository the project-scope skill is installed into \\
\texttt{\mbox{-\allowbreak{}}\mbox{-\allowbreak{}}dest} & \texttt{DEST} & optional & explicit skills folder (overrides \mbox{-}\mbox{-}agent/\mbox{-}\mbox{-}scope) \\
\texttt{\mbox{-\allowbreak{}}\mbox{-\allowbreak{}}force} & -- & optional & replace an existing installation \\
\texttt{\mbox{-\allowbreak{}}\mbox{-\allowbreak{}}dry-\allowbreak{}run} & -- & optional & print the target folder, change nothing \\
\bottomrule
\end{longtable}}

\paragraph{Examples}\leavevmode
\begin{Verbatim}[fontsize=\footnotesize,breaklines,breakanywhere]
scripts/watchover install-skill --agent codex --project ../akeyless-governance
scripts/watchover install-skill --agent claude --scope user
scripts/watchover install-skill --agent copilot --project . --dry-run
scripts/watchover install-skill --agent custom --dest ~/my-agent/skills --force
\end{Verbatim}
\paragraph{Sample run}\leavevmode
% GENERATED by docs/manual/build.py - do not edit; change the source it names instead.
% source: docs/manual/build.py EXAMPLES (executed in a sandbox copy of the repository)
\begin{terminal}
\begin{Verbatim}[fontsize=\scriptsize,breaklines,breakanywhere,commandchars=\\\{\}]
\textcolor{wo-prompt}{\$} \textbf{mkdir -p ../governance-repo && scripts/watchover install-skill --agent codex --project ../governance-repo}
installed akeyless-watchover for OpenAI Codex (CLI, IDE, app) in <tmp>/governance-repo/.agents/skills/akeyless-watchover
\end{Verbatim}
\exitcode{0}
\end{terminal}



\section{\texttt{validate-skill}}\label{kit:validate-skill}
\begin{woentry}
\refrow{Synopsis}{\texttt{scripts/\allowbreak{}watchover validate-\allowbreak{}skill [options]}\hfill\emph{offline}}
\refrow{Does}{Check a skill folder against the Agent Skills specification.}
\refrow{How}{\begin{itemize}[nosep,leftmargin=1.2em]
\item Checks SKILL.md: only specification fields; name 1-64 chars, lowercase/digits/hyphens, equal to the folder; description \textless{}= 1024; compatibility \textless{}= 500; metadata string -\textgreater{} string; body \textless{}= 500 lines; every references/*.md it mentions exists.
\end{itemize}
}
\refrow{Exit codes}{0 valid; 1 problems listed}
\end{woentry}
{\small
\begin{longtable}{@{}>{\raggedright\arraybackslash}p{0.27\linewidth}>{\raggedright\arraybackslash}p{0.17\linewidth}>{\raggedright\arraybackslash}p{0.13\linewidth}>{\raggedright\arraybackslash}p{0.35\linewidth}@{}}
\toprule
\textbf{Option} & \textbf{Value} & \textbf{Required / default} & \textbf{Description} \\ \midrule\endhead
\texttt{path} & \texttt{PATH} & \textbf{required} & skill folder (default: this skill) \\
\bottomrule
\end{longtable}}

\paragraph{Examples}\leavevmode
\begin{Verbatim}[fontsize=\footnotesize,breaklines,breakanywhere]
scripts/watchover validate-skill
scripts/watchover validate-skill dist/akeyless-watchover
\end{Verbatim}
\paragraph{Sample run}\leavevmode
% GENERATED by docs/manual/build.py - do not edit; change the source it names instead.
% source: docs/manual/build.py EXAMPLES (executed in a sandbox copy of the repository)
The skill source:
\begin{terminal}
\begin{Verbatim}[fontsize=\scriptsize,breaklines,breakanywhere,commandchars=\\\{\}]
\textcolor{wo-prompt}{\$} \textbf{scripts/watchover validate-skill}
ok    skill/akeyless-watchover follows the Agent Skills specification
\end{Verbatim}
\exitcode{0}
\end{terminal}

A copy with an invalid name (exit code 1): (The copy is created before and removed after the command shown.)
\begin{terminal}
\begin{Verbatim}[fontsize=\scriptsize,breaklines,breakanywhere,commandchars=\\\{\}]
\textcolor{wo-prompt}{\$} \textbf{scripts/watchover validate-skill /tmp/wo-broken      # name: Akeyless_Watchover, folder wo-broken}
FAIL  name: 1-64 characters, lowercase letters, digits and single hyphens, no leading/trailing hyphen
\end{Verbatim}
\exitcode{1}
\end{terminal}



\section{\texttt{mcp}}\label{kit:mcp}
\begin{woentry}
\refrow{Synopsis}{\texttt{scripts/\allowbreak{}watchover mcp [options]}\hfill\emph{offline}}
\refrow{Does}{Serve the kit, the docs and the skill commands to any MCP client (stdio).}
\refrow{How}{\begin{itemize}[nosep,leftmargin=1.2em]
\item JSON-RPC 2.0 over stdin/stdout. Tools are generated from each command's options and run \texttt{scripts/\allowbreak{}watchover \textless{}command\textgreater{}} in the kit root; output and exit code are returned.
\item Resources: SKILL.md, docs, model files. Prompts: one per skill command (skill instructions + request).
\item Commands that call Akeyless are listed only with \mbox{-}\mbox{-}live.
\end{itemize}
}
\refrow{Exit codes}{runs until the client closes stdin}
\refrow{Environment}{\texttt{WATCHOVER\_\allowbreak{}MCP\_\allowbreak{}LIVE}: 1 = also expose the commands that call Akeyless (verify, export, account-settings, token)\newline \texttt{WATCHOVER\_\allowbreak{}REPO}: kit root the tools run in (default: this repository)}
\end{woentry}
{\small
\begin{longtable}{@{}>{\raggedright\arraybackslash}p{0.27\linewidth}>{\raggedright\arraybackslash}p{0.17\linewidth}>{\raggedright\arraybackslash}p{0.13\linewidth}>{\raggedright\arraybackslash}p{0.35\linewidth}@{}}
\toprule
\textbf{Option} & \textbf{Value} & \textbf{Required / default} & \textbf{Description} \\ \midrule\endhead
\texttt{\mbox{-\allowbreak{}}\mbox{-\allowbreak{}}repo} & \texttt{REPO} & optional & kit root the tools run in and the model/docs are read from (default: \$WATCHOVER\_REPO or this repository) \\
\texttt{\mbox{-\allowbreak{}}\mbox{-\allowbreak{}}live} & -- & optional & also expose the commands that call Akeyless (same as WATCHOVER\_MCP\_LIVE=1) \\
\texttt{\mbox{-\allowbreak{}}\mbox{-\allowbreak{}}list} & -- & optional & print the tools, resources and prompts, then exit \\
\bottomrule
\end{longtable}}

\paragraph{Examples}\leavevmode
\begin{Verbatim}[fontsize=\footnotesize,breaklines,breakanywhere]
scripts/watchover mcp
WATCHOVER_MCP_LIVE=1 scripts/watchover mcp --repo ~/src/akeyless-governance
\end{Verbatim}
\paragraph{Sample run}\leavevmode
% GENERATED by docs/manual/build.py - do not edit; change the source it names instead.
% source: docs/manual/build.py EXAMPLES (executed in a sandbox copy of the repository)
What an MCP client will see (offline commands only):
\begin{terminal}
\begin{Verbatim}[fontsize=\scriptsize,breaklines,breakanywhere,commandchars=\\\{\}]
\textcolor{wo-prompt}{\$} \textbf{scripts/watchover mcp --list}
tools:     watchover_help, watchover_test, watchover_simulate, watchover_forwarder, watchover_conjur, watchover_build_skill, watchover_validate_skill
resources: watchover://skill/SKILL.md, watchover://docs/README.md, watchover://docs/agents.md, watchover://docs/api-vs-cli.md, watchover://docs/architecture.md, watchover://docs/brownfield-import.md, watchover://docs/cli.md, watchover://docs/conjur-migration.md, watchover://docs/limitations.md, watchover://docs/model-reference.md, watchover://docs/path-and-role-design.md, watchover://docs/pipeline.md, watchover://docs/pki.md, watchover://docs/rbac-model.md, watchover://docs/terraform.md, watchover://docs/testing-and-verification.md, watchover://docs/usage-guide.md, watchover://model/access_expectations.yaml, watchover://model/auth_methods.yaml, watchover://model/folders.yaml, watchover://model/gateway.yaml, watchover://model/governance.yaml, watchover://model/pki.yaml, watchover://model/roles/payments.yaml, watchover://model/roles/platform.yaml
prompts:   watchover-help, watchover-design, watchover-change, watchover-review, watchover-pipeline, watchover-import, watchover-conjur, watchover-pki, watchover-harden, watchover-gateway, watchover-verify, watchover-cli, watchover-explain
\end{Verbatim}
\exitcode{0}
\end{terminal}

With the live commands enabled:
\begin{terminal}
\begin{Verbatim}[fontsize=\scriptsize,breaklines,breakanywhere,commandchars=\\\{\}]
\textcolor{wo-prompt}{\$} \textbf{scripts/watchover mcp --list --live | head -n 1}
tools:     watchover_help, watchover_test, watchover_simulate, watchover_verify, watchover_export, watchover_forwarder, watchover_conjur, watchover_account_settings, watchover_token, watchover_build_skill, watchover_validate_skill
\end{Verbatim}
\exitcode{0}
\end{terminal}

A raw MCP session over stdio (what a client sends), output shortened per message: (mcp\_demo.py is the short Python script embedded in docs/manual/build.py.)
\begin{terminal}
\begin{Verbatim}[fontsize=\scriptsize,breaklines,breakanywhere,commandchars=\\\{\}]
\textcolor{wo-prompt}{\$} \textbf{python3 mcp_demo.py   # sends initialize, tools/call watchover_simulate and prompts/get watchover-change to `scripts/watchover mcp` over stdio and prints a summary of each answer}
initialize -> \{'name': 'watchover-mcp', 'title': 'akeyless-watchover', 'version': '1.5.0'\} protocol 2025-06-18
tools/call watchover_simulate -> isError False
  $ scripts/watchover simulate --model terraform/model --expectations terraform/model/access_expectations.yaml
  (exit code 0)
  ... 0 mismatch(es)
prompts/get watchover-change -> 19264 characters:
  Follow the akeyless-watchover skill instructions below. Run its `change` command with: give fraud-analysts rea...
\end{Verbatim}
\exitcode{0}
\end{terminal}



\section{\texttt{help} and \texttt{version}}\label{kit:help}
\begin{woentry}
\refrow{Synopsis}{\texttt{scripts/watchover help [topic] [--markdown]}; \texttt{scripts/watchover version}\hfill\emph{offline}}
\refrow{Does}{Prints the overview, one command's entry, the skill commands, or this reference as Markdown; \texttt{version} prints the version. No argument = \texttt{help}.}
\end{woentry}
{\small
\begin{longtable}{@{}>{\raggedright\arraybackslash}p{0.27\linewidth}>{\raggedright\arraybackslash}p{0.17\linewidth}>{\raggedright\arraybackslash}p{0.13\linewidth}>{\raggedright\arraybackslash}p{0.35\linewidth}@{}}
\toprule
\textbf{Option} & \textbf{Value} & \textbf{Required / default} & \textbf{Description} \\ \midrule\endhead
\texttt{topic} & \texttt{TOPIC} & \textbf{required} & a kit command name or 'skill'; omitted = overview \\
\texttt{\mbox{-\allowbreak{}}\mbox{-\allowbreak{}}markdown} & -- & optional & print this reference as Markdown (docs/usage-guide.md) \\
\bottomrule
\end{longtable}}

\paragraph{Sample run}\leavevmode
% GENERATED by docs/manual/build.py - do not edit; change the source it names instead.
% source: docs/manual/build.py EXAMPLES (executed in a sandbox copy of the repository)
\begin{terminal}
\begin{Verbatim}[fontsize=\scriptsize,breaklines,breakanywhere,commandchars=\\\{\}]
\textcolor{wo-prompt}{\$} \textbf{scripts/watchover help}
watchover 1.5.1 - akeyless-watchover kit

Usage: scripts/watchover <command> [options]

Commands:
  test              Run every offline check (what CI runs).
  simulate          Offline: evaluate access_expectations.yaml against the model.
  verify            Live: compare expectations with describe-permissions run as each persona's CLI profile.  [needs Akeyless]
  export            Export existing config and translate it: YAML model + import blocks, or plain Terraform (--hcl).  [needs Akeyless]
  forwarder         Scaffold a new event forwarder (model entry or HCL); --import-existing adopts one by name.
  conjur            Translate Conjur policy into model roles + identity bindings.
  account-settings  Apply account_settings_cli (settings the provider cannot manage) with the CLI.  [needs Akeyless]
[... 12 lines omitted ...]
\end{Verbatim}
\exitcode{0}
\end{terminal}



